\documentclass[conference]{IEEEtran}
\usepackage[utf8]{inputenc}
\usepackage[T1]{fontenc}
\usepackage{cite}
\usepackage{amsmath,amssymb}
\usepackage{graphicx}
\usepackage{booktabs}
\usepackage{url}

\title{A Preregistered Audit of Dataset Contamination Controls in LLM-for-NetOps Benchmarks}

\author{
\IEEEauthorblockN{Autonomous AI Researcher}
\IEEEauthorblockA{CNSM 2026 Agentic AI Researcher Track}
}

\begin{document}

\maketitle

\begin{abstract}
We preregistered and executed a paired, confirmatory audit to test whether conservative deterministic contamination detectors (exact‑match + fuzzy 5‑gram + provenance heuristics) flag items in a reproducible 150‑item NetOps sample and whether applying preregistered contamination controls changes validator‑based success rates. Crucially, the executed sample used provenance‑stable synthetic tasks (source\_identifier prefix "synthetic-netops:") and shared generation semantics (independent\_condition\_generation = false), recorded in the archived model and task manifests; these two execution facts materially limited the audit's sensitivity to generation‑level contamination. No preregistered high‑confidence contamination flags occurred and therefore no guarded exclusions or replacements were performed. All 150 baseline and matched guarded episodes passed the deterministic validator (baseline = 150/150, guarded = 150/150). The paired success‑rate difference = 0.0 percentage points (95\% CI [0.0, 0.0]); exact McNemar two‑sided p = 1.0. We reconcile the preregistration\ensuremath{\rightarrow}execution gap with artifact‑grounded diagnosis, explain why the run is degenerate for detecting public‑benchmark leakage, and provide concrete, artifact‑anchored recommendations for audit designs that would be sensitive to contamination in web‑sourced benchmarks.
\end{abstract}

\section{Introduction}
Motivation. Large language models (LLMs) are increasingly applied to NetOps tasks such as intent‑to‑configuration translation and automated configuration repair. Operational decisions based on benchmark results require that measured performance reflect model capability rather than dataset contamination or templated item leakage into model training. Meta‑analyses and recent surveys call for contamination‑aware evaluations and validator‑backed oracles to avoid inflated reliability claims.

Research question and contribution. We preregistered and executed an artifact‑anchored audit to answer: under conservative deterministic contamination detectors (exact normalized token equality; fuzzy 5‑gram shingle overlap with preregistered thresholds; provenance timestamp heuristics) and a guarded exclusion policy, how many high‑confidence flags arise in a reproducible 150‑item NetOps sample, and how much does applying those controls change a single hosted LLM's validator pass rate? Our contributions are: (1) a fully preregistered confirmatory experiment (study\_id = contamination-audit-netops-2026-v1) executed end‑to‑end and archived; (2) exact, artifact‑verified confirmatory statistics on Npairs = 150 with deterministic validator traces; (3) an explicit preregistration\ensuremath{\rightarrow}execution reconciliation that identifies why the run produced a degenerate null and prescribes artifact‑grounded design changes to increase audit sensitivity.

Scope and bounded claims. The audit intentionally targeted a methodological question about preregistered contamination controls, not an empirical prevalence estimate for public web‑sourced benchmarks. The executed confirmatory sample used provenance‑stable synthetic tasks (source\_identifier prefix "synthetic-netops:") produced by the deterministic task generator and employed shared\_initial\_candidate generation semantics (independent\_condition\_generation = false). Because both facts are recorded in the archived model and task manifests, they limit inference: the observed zero‑flag, zero‑difference outcome applies to this synthetic sampling and generation regime and does not imply absence of contamination in independent public benchmarks.

\section{Related Work}
LLM‑driven NetOps systems and formal validators. Recent system and benchmark work documents the feasibility of translating intent into device configuration artifacts and of using LLMs as repair generators in controlled settings. Intent‑driven configuration papers demonstrate that LLMs can produce syntactically plausible and often functionally correct configurations when tasks are constrained and prompts are carefully designed [7]. Benchmark and systems studies focused on configuration repair report that generator output quality improves when paired with deterministic validators or iterative validate--repair loops; these generate‑then‑verify patterns are an increasingly accepted design for raising practical reliability [1,3]. Complementary to these LLM‑centric evaluations, classical network verification research (e.g., header‑space analyses and localized subspecification approaches) supplies scalable, property‑focused oracles for reachability and policy checks; these formal tools are practical building blocks for generate--validate pipelines that check functional correctness beyond surface syntactic match [4].

Agentic and integrated repair architectures. More recent work evaluates agentic or tool‑augmented pipelines that combine LLM generation, retrieval/context, and external verifiers. Empirical comparisons show that agentic integrations with explicit verifier calls and iterative repairs typically reduce regressions relative to naive single‑shot generation, but they do not eliminate per‑case failures and their effectiveness depends on the verifier's coverage and the agent's tool contract [2]. System papers that emphasize tighter repair chains (localize\ensuremath{\rightarrow}propose\ensuremath{\rightarrow}validate\ensuremath{\rightarrow}refine) report higher average repair rates on curated incident datasets, yet also highlight sensitivity to task heterogeneity, topology realism, and the exact verifier semantics used in evaluation [3]. These findings motivate using deterministic validators as the primary binary oracle in contamination‑aware audits, since validators operationalize functional pass/fail outcomes that matter to network operators.

Evaluation‑validity critiques and contamination detection methods. Independent surveys and conceptual frameworks have raised two recurring risks for LLM evaluations relevant to NetOps: dataset contamination (items or paraphrases appearing in model training data) and construct invalidity (benchmarks dominated by low‑complexity, templated prompts that overstate operational readiness) [8,6,5]. The literature recommends layered defenses: deterministic exact‑match fingerprinting, fuzzy n‑gram (shingle) overlap detectors with sensitivity analyses over multiple thresholds, provenance timestamp checks versus model training cutoffs, and human adjudication for borderline cases. These recommended practices inform our preregistered detector suite (exact canonical equality + fuzzy 5‑gram overlap thresholds + provenance heuristics) and our guarded exclusion policy. Prior work also stresses preserving detector outputs (full overlap score distributions) to enable post‑hoc assessment of threshold choices; this practical point motivated our archival design, and it underpins our recommendation to persist per‑item overlap statistics even when no flags cross the preregistered thresholds [8].

Positioning relative to prior studies. Unlike broad capability benchmarks that report absolute performance across many models, our study is a preregistered, artifact‑anchored audit that asks whether conservative, preregistered contamination controls would change validator‑measured success on a reproducible sample under explicitly recorded execution semantics. Prior benchmark and system work assessed LLM repair efficacy and proposed integrated agentic architectures; contamination‑awareness and detector design have been discussed in surveys and conceptual proposals but have seen fewer fully preregistered, verifier‑backed confirmatory audits. Our contribution complements those lines by (a) operationalizing a conservative detector suite in a preregistration; (b) using a deterministic verifier as the functional oracle; and (c) providing an artifact‑grounded reconciliation when preregistered choices (synthetic provenance + shared generation semantics) rendered the confirmatory paired test degenerate. This positioning clarifies that our findings diagnose audit design sensitivity rather than claiming empirical prevalence of contamination in independent web‑sourced benchmarks.

\section{Methodology}
Preregistration and primary estimand. The experiment was preregistered (study\_id = contamination-audit-netops-2026-v1) and froze analysis and exclusion rules before execution. The primary estimand is the paired success‑rate difference (guarded - baseline) across item\ensuremath{\times}model pairs, tested by exact McNemar two‑sided test and summarized with a paired bootstrap percentile 95\% CI (10,000 resamples, seed = 1729) as specified in the preregistered paired\_binary\_analysis\_v1 executor.

Detectors and guarded policy (preregistered). The preregistered detectors were: (1) exact normalized token equality after canonicalization; (2) fuzzy 5‑gram shingle overlap (Jaccard/overlap) with preregistered high‑confidence threshold >= 0.80 and sensitivity thresholds at 0.65 and 0.50; and (3) provenance timestamp heuristics comparing item timestamps to model cutoff. The guarded policy deterministically excludes only high‑confidence flagged items (exact match OR fuzzy overlap >= 0.80) and replaces them via deterministic retemplating/topology augmentation from the same generator with fixed seeds; replacements must meet an embedding similarity constraint per the preregistration.

Sampling, generation semantics, and the critical execution choice. The preregistration allowed deterministic task generators for reproducibility and permitted shared generation semantics when appropriate. The executed run used deterministic\_netops\_task\_generator\_v5 to produce synthetic tasks (source\_identifier prefix "synthetic-netops:") for the confirmatory sample. Execution settings used shared\_initial\_candidate semantics (independent\_condition\_generation = false): when no exclusions occur, one shared model generation per item is reused for both baseline and guarded episodes. These execution settings are recorded in the archived model\_configuration and task\_manifest artifacts. The preregistered contract also planned replacement rules, sensitivity analyses, and a maximum of 300 model calls (one per condition per item) if independent generations were used.

Verifier and scoring rule. The verification oracle is deterministic\_netops\_validator\_v5. For intent\ensuremath{\rightarrow}configuration tasks the validator applies parse\ensuremath{\rightarrow}state‑update\ensuremath{\rightarrow}property checks and returns a binary pass/fail along with detailed traces (observed\_touched\_field\_count, parsed\_command\_count, required\_command\_count, violation codes). The preregistered scoring rule is full\_intent\_constraint\_validation: outputs must achieve the intended final state and respect workflow constraints. Validator traces are archived for every episode and provide the empirical basis for the binary outcomes used in the paired analysis.

Planned sensitivity and power. The preregistration included sensitivity analyses: treating failed model calls as failures; evaluating fuzzy thresholds at 0.50, 0.65, 0.80; and subgrouping by templated vs realistic items. The power guidance documented that a multi‑model or two‑call per item design (300 independent generations) would detect smaller paired differences (\ensuremath{\approx}4--6 percentage points under specified assumptions), while the single‑model single‑call shared‑generation regime has reduced sensitivity. These planning notes are part of the preregistration archive and informed our interpretation of the executed null result.

\section{Results}
Execution accounting and reconciliation. The archived execution manifest records planned\_episode\_count = 300 (150 items \ensuremath{\times} 2 conditions) and completed\_episode\_count = 300 with failed\_episode\_count = 0. Because the run used shared\_initial\_candidate semantics and no high‑confidence contamination exclusions occurred, the pipeline consumed model\_calls\_used = 150 (one shared model generation per item) rather than two independent generations per condition; provider call auditing corroborates hosted\_model\_call\_event\_count = 150, cache\_miss\_count = 150 and cross\_condition\_cache\_key\_reuse\_observed = false. The deterministic\_reconciliation artifact confirms complete\_pair\_count = 150 and internal accounting consistency (n\_11 = 150, n\_10 = n\_01 = n\_00 = 0).

Contamination detectors and archived evidence of absence. Under the preregistered high‑confidence rule (exact match OR fuzzy 5‑gram overlap >= 0.80) the detector workflow produced zero high‑confidence flags for the executed sample. The analysis artifact contamination\_summary.csv is empty (no flagged rows), and the run therefore performed no deterministic exclusions or replacements. Per the preregistered detector implementation, detailed per‑item fuzzy overlap scores and replacement traces are persisted only when threshold crossings or replacements occur; because no items met the high‑confidence threshold, the archive contains the detector's flag events (none) but not a full per‑item overlap table. We therefore report the artifact‑grounded null finding (no high‑confidence flags) and emphasize that the absence of persisted overlap summaries is itself an execution artifact with diagnostic meaning (see below).

Validator outcomes: concordant passes and formal statistics. The deterministic\_netops\_validator\_v5 returned binary pass for every baseline episode (baseline\_success\_count = 150/150, baseline\_success\_rate = 1.0) and for every guarded episode (guarded\_success\_count = 150/150, guarded\_success\_rate = 1.0). The paired contingency table is n11 = 150, n10 = 0, n01 = 0, n00 = 0. The primary estimand (paired\_success\_rate\_difference\_guarded\_minus\_baseline) = 0.0 percentage points. The exact McNemar two‑sided test on these counts yields p = 1.0. The preregistered paired bootstrap percentile CI (10,000 resamples, seed = 1729) is [0.0, 0.0]. The zero‑width CI and p = 1.0 are a direct consequence of the absence of any discordant pairs; these are artifact‑verified outcomes derived from the archived paired\_contingency\_table.csv and condition\_summary.csv.

Representative diagnostic episodes. Representative archived episode traces illustrate validator behavior and the nature of shared generations. For example, pair‑000001 (task‑000001) has normalized\_response identical to the reference, parsed\_command\_count = 4 and required\_command\_count = 4, observed\_touched\_field\_count = 3, and validator verdict = VALID\_CONFIGURATION. Pair‑000002 (task‑000002) shows parsed\_command\_count = 2 (required\_command\_count = 2, level = hard) and validator verdict = VALID\_CONFIGURATION. Pair‑000003 (task‑000003) shows parsed\_command\_count = 6 (required\_command\_count = 6, level = hard) and validator verdict = VALID\_CONFIGURATION. These exemplar traces (archived per‑episode validator\_trace records) demonstrate that when the model output aligns with the preregistered intent and workflow constraints the validator produces a detailed, machine‑readable confirmation; in this execution these per‑episode confirmations are uniformly positive.

Why the paired test was uninformative for generation‑level contamination. Two artifact‑recorded execution facts explain the degenerate concordance. First, the confirmatory sample comprised provenance‑stable, deterministically generated synthetic items (source\_identifier prefix "synthetic-netops:"), so exact canonicalized matches and very high fuzzy 5‑gram overlaps to public corpus fingerprints were unlikely a priori. Second, independent\_condition\_generation = false (shared\_initial\_candidate) was in effect: when no high‑confidence flags triggered exclusions the same single model output was reused for both baseline and guarded conditions. Together these facts guarantee that, absent deterministic exclusions, baseline and guarded episodes for each item are identical in content and validator outcome---producing n11 = 150 and zero discordance by construction. The execution manifest and model\_configuration artifact explicitly encode these settings and thus ground this diagnosis in archived evidence rather than post hoc speculation.

Detector persistence policy and the missing overlap distribution. The preregistered detector persisted detailed fuzzy overlap scores only when an item crossed a threshold or when replacements were enacted; the archived absence of persisted per‑item overlap scores therefore reflects the detector's control flow under the executed conditions (no threshold crossings). This design choice reduced the post‑hoc ability to report summary overlap statistics (min/median/percentiles) for the executed sample from the archive. The missing overlap distribution is an execution artifact: it indicates that no items came close enough to the high‑confidence threshold to invoke persistence logic, but it does not reveal whether many items would have had moderate overlaps below the threshold (the archive contains no full table of low‑confidence overlaps). This factual gap motivates our concrete recommendation to persist per‑item overlap scores regardless of threshold crossings in future audits to enable transparent reporting of the overlap distribution and sensitivity analyses.

Sensitivity to preregistered design choices. The preregistration explicitly documented alternative, higher‑sensitivity designs: (a) independent\_condition\_generation = true so baseline and guarded generations are independent (doubling model calls to 300) and allowing per‑item discordance to appear even without deterministic flagging; (b) multi‑model replication to reveal model‑specific memorization signals; and (c) lower fuzzy thresholds or paraphrase‑robust detectors plus human adjudication. The archived power guidance in the preregistration explains that a 300 independent‑call design or a multi‑model plan would materially lower the detectable paired difference (planned target \ensuremath{\approx}5 percentage points at 80\% power under conservative assumptions). Because the executed run used the lower‑sensitivity shared‑generation regime and synthetic sampling, the artifact‑grounded conclusion is limited to this regime: deterministic detectors did not trigger exclusions and validators returned universal passes for the shared outputs.

Summary of archival provenance supporting the numerical results. All numeric summaries above---completed\_episode\_count = 300, model\_calls\_used = 150, baseline\_success\_count = 150, guarded\_success\_count = 150, n11 = 150, paired bootstrap CI [0.0,0.0] (10,000 resamples, seed = 1729), and McNemar p = 1.0---are computed by the preregistered analysis executor and are recorded in analysis/results.json and in the small CSV artifacts condition\_summary.csv and paired\_contingency\_table.csv. The contamination\_summary.csv artifact contains no flagged rows for this execution. Per‑episode raw responses and validator\_trace objects are archived for independent verification of every binary outcome used in the paired analysis.

\section{Discussion}
Why the result is degenerate: artifact‑grounded diagnosis. Two archived execution facts together explain the degenerate null outcome. First, the confirmatory sample was drawn entirely from provenance‑stable synthetic tasks generated by deterministic\_netops\_task\_generator\_v5 (task entries carry the "synthetic‑netops:" source\_identifier). Synthetic provenance makes exact normalized token equality and high fuzzy 5‑gram overlap with public corpus fingerprints unlikely by construction, and therefore the preregistered high‑confidence detector rule (exact match OR 5‑gram overlap >= 0.80) had little opportunity to trigger. Second, execution semantics set independent\_condition\_generation = false (shared\_initial\_candidate). When the detector produced no high‑confidence exclusions, the pipeline reused a single model generation for both baseline and guarded episodes; this reuse is recorded in the model and provider accounting artifacts and reduces independent variance between conditions. Both facts are present in the archived execution/model\_configuration.json and execution/task\_manifest.jsonl artifacts and confirmed by the deterministic\_reconciliation outputs. Together these facts deterministically produce identical baseline/guarded outputs per item and therefore make the paired test uninformative for generation‑level contamination effects in this run.

Mechanics of the detector workflow and reporting consequences. The preregistered detector was implemented to persist detailed per‑item overlap scores only when threshold crossings or replacement handling were invoked (to limit storage and to focus archival attention on resolved flags and replacements). Because no items crossed the high‑confidence threshold, the detector produced no persisted per‑item overlap table in analysis/contamination\_summary.csv. This implementation choice---record only flagged events---yields a clear archival signal when contamination is detected but produces a blind spot for audits that later wish to inspect the full overlap distribution. The artifact evidence therefore shows an absence of high‑confidence flags but does not let us characterize how close items were to the threshold in aggregate. Practically, audit pipelines should persist per‑item detector scores at least as summary quantiles so a null flag set can still be interrogated post hoc; we recommend this change precisely because the current archive lacks the overlap distribution needed to rule out near‑threshold cases.

Effect on estimand sensitivity and power. The preregistration explicitly contrasted two plausible execution regimes: (A) independent condition generation (two independent model calls per item, one for baseline and one for guarded) and (B) shared generation (single call reused across conditions when exclusions do not apply). Independent generation increases the number of independent model outputs and therefore raises the probability of observing cross‑condition discordance arising from generation‑level contamination, paraphrase variance, or randomness. The executed shared‑generation regime collapses those independent draws into one and therefore reduces observable discordance unless the detector deterministically excludes an item. The execution manifest shows model\_calls\_used = 150 versus a planned maximum of 300; this arithmetic difference is the core practical reason the paired estimator had no discordant pairs. The preregistration's power guidance described this tradeoff qualitatively; the artifacts make the tradeoff concrete for this run and explain why the paired test could not detect generation‑level leakage under the chosen semantics.

Validator coverage and semantic failure modes. The deterministic\_netops\_validator\_v5 is a practical, deterministic oracle for the operational properties it encodes: parse\ensuremath{\rightarrow}state updates, required command counts, dependency rules, and final‑state checks. Validator traces archived per episode (validation\_before/validation\_after) demonstrate that the validator detected and recorded command counts, touched fields, and violation codes for every candidate. However, artifact evidence also makes clear the validator's scope is constrained to the preregistered intent constraints and modeled properties; it does not claim coverage of unmodeled semantic failures (for example, subtle ordering effects outside the specified workflow\_policies, vendor CLI idiosyncrasies not represented in the generator, or emergent cross‑task interactions). Consequently, a passed validation trace is necessary evidence of meeting the modeled intent but not sufficient evidence of comprehensive production safety. This distinction is present in the archived validator\_trace fields and motivates our recommendations for richer oracle coverage and human adjudication for borderline or high‑risk items.

Concrete reproducibility and artifact inspection guidance. The execution bundle contains the artifacts that substantiate the diagnosis: execution/model\_configuration.json documents the requested execution semantics (declared\_model\_version = "gpt-5-mini", independent\_condition\_generation setting), execution/task\_manifest.jsonl records each sampled task with source\_identifier and generation\_seed, analysis/contamination\_summary.csv is empty (no high‑confidence flags), analysis/paired\_contingency\_table.csv records n11 = 150 and zero discordant pairs, and analysis/paired\_difference\_figure.svg contains the bootstrap distribution used to compute the CI. Inspecting these artifacts in the archive allows an auditor to verify the exact pipeline behavior without recomputing model outputs. We highlight that persistent per‑item detector scores were intentionally omitted when flags were absent; auditors who require those scores should request pipelines that persist per‑item summaries regardless of flag status.

Design prescriptions to increase audit sensitivity. Artifact‑grounded lessons suggest specific, implementable changes: (1) sample provenance‑rich, web‑sourced items (with stable URLs/DOIs and timestamps) to exercise the fuzzy‑overlap and provenance heuristics; (2) set independent\_condition\_generation = true and execute two independent model calls per item to create the independent comparisons that the paired test is designed to detect; (3) increase model diversity (multi‑model replication) to surface model‑specific memorization patterns; (4) persist per‑item fuzzy‑overlap scores and canonicalized fingerprints even when thresholds are not crossed so the overlap distribution is reconstructible; and (5) add paraphrase‑robust detectors (semantic similarity beyond n‑gram overlap) plus human adjudication for borderline results. Each prescription follows directly from the execution artifacts: the model\_calls\_used shortfall, the empty contamination\_summary, and the recorded shared\_initial\_candidate semantics.

Operational implications and auditor tradeoffs. Auditors face tradeoffs between reproducibility and sensitivity. Deterministic synthetic sampling and shared generation maximize reproducibility and reduce compute, but those same choices can eliminate the observable signal needed to detect training‑data leakage at generation time. Conversely, independent generations and multi‑model designs increase sensitivity but require more calls, storage, and non‑determinism management. The archived execution demonstrates that reproducibility‑focused choices can produce a valid confirmatory test that nevertheless answers a narrower methodological question than originally targeted; explicit preregistration of those tradeoffs (as in our archive) is therefore essential for correct interpretation.

Broader methodological takeaways. Artifact‑anchored preregistration remains a powerful practice: freezing detectors, thresholds, replacement rules, and analysis executors before observing outcomes prevents post hoc rule changes. However, the audit shows that preregistration must also commit to sampling and generation semantics that align with the estimand's sensitivity requirements. In short, auditors should preregister not only detectors and thresholds but also the independence structure of model draws and the persistence policy for detector outputs; the execution artifacts here show all three choices materially affected inferential ability.

\section{Limitations}
\begin{itemize}
\item Single‑model execution (gpt‑5-mini) --- results do not generalize across other LLM families, sizes, or training corpora.
\item Sampled items were provenance‑stable synthetic/deterministically generated tasks (source\_identifier prefix "synthetic-netops:") for this run; absence of high‑confidence flags here may not reflect contamination prevalence in real‑world, web‑crawled benchmarks.
\item Generation semantics used shared\_initial\_candidate across baseline and guarded conditions (independent\_condition\_generation = false), producing identical model outputs when no exclusions occurred and thus limiting sensitivity to interventions that rely on independent re‑generation.
\item Contamination detection relied on preregistered conservative heuristics (exact match and fuzzy 5‑gram overlap thresholds); subtler contamination patterns (paraphrased memorization, partial leakage) may escape these heuristics and require richer provenance or human adjudication.
\item Passing the deterministic\_netops\_validator\_v5 indicates satisfaction of the checked functional constraints but does not imply comprehensive production readiness or absence of semantic regressions outside the validator's scope.
\end{itemize}

\section{Conclusion}
We executed a fully preregistered, artifact‑anchored paired audit of conservative contamination detectors for LLM‑for‑NetOps evaluation. In the reproducible 150‑pair synthetic sample we observed zero preregistered high‑confidence contamination flags and identical validator pass rates (150/150 baseline, 150/150 guarded). The degenerate null outcome is artifact‑supported and stems from provenance‑stable synthetic sampling and shared generation semantics; it does not imply absence of contamination in web‑sourced benchmarks or failure of the preregistered detectors in different sampling regimes. We provide an explicit preregistration\ensuremath{\rightarrow}execution reconciliation, confirm deterministic reconciliation checks passed, and recommend concrete design changes---provenance‑rich sampling, independent generation per condition, multi‑model scope, paraphrase‑sensitive detectors, and matched power calculations---to increase audit sensitivity in future studies. All experiment artifacts, validator traces, analysis inputs/outputs, and execution manifests are archived and available as recorded in the Disclosure Statement below.

\section*{Disclosure Statement}
This study was executed under the autonomous‑run preregistration contamination-audit-netops-2026-v1 (preregistration SHA‑256 recorded in the archived bundle). LLM used: gpt-5-mini (declared\_model\_version = "gpt-5-mini"). Orchestration adapter: hosted\_netops\_gvr\_v1. Deterministic task generator: deterministic\_netops\_task\_generator\_v5 (version 5.0). Deterministic validator: deterministic\_netops\_validator\_v5 (version 5.0). Representative executed artifacts (by path) include execution/model\_configuration.json, execution/task\_manifest.jsonl, execution/raw\_results.jsonl, analysis/paired\_contingency\_table.csv, and analysis/contamination\_summary.csv; the full execution bundle contains raw responses, validator traces, execution logs, and the transformation manifest. Exact immutable initial master‑prompt reference: provenance/master\_prompt.txt (SHA‑256: 1872df1e\allowbreak{}1805d2d9\allowbreak{}6940456c\allowbreak{}a016bd66\allowbreak{}5d1d5196\allowbreak{}add77f5a\allowbreak{}cdf1582b\allowbreak{}b39b15ba). Topic constraints (Generative AI and LLMs for NetOps) were selected by the human Research Director prior to the autonomy lock. After the autonomy lock, the AI system autonomously performed literature search, gap selection, preregistration, experiment execution, analysis, manuscript drafting, peer‑review simulation, and revision. All generation prompts, model outputs, validator traces, sampling seeds, replacement rules, execution logs, and analysis artifacts were produced and archived by the autonomous pipeline and are included in the execution bundle. No manual human editing of technical content, analyses, references, figures, tables, captions, or the Disclosure Statement occurred after the autonomy lock. Pre‑lock development and rehearsal runs occurred (permitted) but pre‑lock artifacts were not used as empirical evidence for this final study unless re‑created under the locked pipeline. The execution followed preregistered missingness, retry, and exclusion policies; no deviations occurred.

\begin{thebibliography}{99}
\bibitem{ref1} https://arxiv.org/abs/2604.22513
\bibitem{ref2} https://doi.org/10.48550/arxiv.2606.06212
\bibitem{ref3} https://doi.org/10.48550/arxiv.2605.22092
\bibitem{ref4} http://citeseerx.ist.psu.edu/viewdoc/summary?doi=10.1.1.300.8659
\bibitem{ref5} https://doi.org/10.1007/s11704-026-60308-3
\bibitem{ref6} Buğra Alperen Uluırmak, Rıfat Kurban, "EvalSafetyGap: A Hybrid Survey and Conceptual Framework for LLM Evaluation-Safety Failures", 2026
\bibitem{ref7} Nguyen Van Tu, Sukhyun Nam, James Won‐Ki Hong, "Intent‐Based Network Configuration Using Large Language Models", 2024, 10.1002/nem.2313
\bibitem{ref8} Lei Huang, Weijiang Yu, Weitao Ma, Weihong Zhong, Zhangyin Feng, Haotian Wang, Qianglong Chen, Weihua Peng, Xiaocheng Feng, Bing Qin, Ting Liu, "A Survey on Hallucination in Large Language Models: Principles, Taxonomy, Challenges, and Open Questions", 2024, 10.1145/3703155
\end{thebibliography}

\end{document}
